\documentclass{article}

\usepackage{amsmath}
\usepackage{amssymb}
\usepackage{geometry}

\geometry{margin=1in}

\begin{document}

\begin{flushright}
CSE 21 Fa26\\
Homework 2\\
Submitted: 10/12/2026
\end{flushright}

\begin{center}
David Washington(200484943), Judah Boyce(200453537), Osinachi Agaranna(200455712), Nick Young(200455733)
\end{center}

\begin{enumerate}
    \item  Consider the set H of all non-negative integer solutions to the equation:
    \[ a_1 + a_2 + a_3 + a_4 + a_5 + a_6 = 30 \]
    \begin{enumerate}
        \item What is the probability that you get the element: (5, 5, 5, 5, 5, 5)? Shoudl we call it like event A so the possibilit that even A occures is equal to the possible outcomes / the total outcomes. Since this is a probability question you need to find the fraction of the possible outcomes over total outcomes.
        \begin{enumerate}
            \item  In this case the total outcomes is: $6^{30}$ ( all 30 dice roll has six potential options).
            \item  The total possible outcomes is the product of the number of ways to choose 5 dice rolls for that number, given the number of dice rolls left:
            \[  \binom{30}{5} \cdot \binom{25}{5} \cdot \binom{20}{5} \cdot \binom{15}{5} \cdot \binom{10}{5} \cdot \binom{5}{5} \] 
        \end{enumerate}
        Solution:
        The possibility that you get the element (5, 5, 5, 5, 5, 5): 
        \[ P(A) = \frac{\binom{30}{5} \cdot \binom{25}{5} \cdot \binom{20}{5} \cdot \binom{15}{5} \cdot \binom{10}{5} \cdot \binom{5}{5}}{ 6^{30}} \] 
        \item What is the probability that at least one of the entries is 0?
        Event A is the probability that at least one of entries is a 0
        So the total probability of all outcomes is again as stated earlier: $6^{30}$. To find the number of ways to have at least one zero entry we use the inclusion exclusion rule starting from choosing oen of the 6 characters to be zero untill all but one was rolled zero times. The representation of the inclusion-exclusion is as follows:
        \[ \binom{6}{1} 5^{30} - \binom{6}{2} 4^{30} + \binom{6}{3} 3^{30} - \binom{6}{4} 2^{30}  + \binom{6}{5} 1^{30} \] 
        we use the inclusion exclusion rule to remove the double counted values (give a short exmaple). 
        \textbf{Solution: The possibility of event A, that at least one of the entris is a 0:}
        \[ P(A) = \frac{\binom{6}{1} 5^{30} - \binom{6}{2} 4^{30} + \binom{6}{3} 3^{30} - \binom{6}{4} 2^{30}  + \binom{6}{5} 1^{30}}{6^{30}} \]   

        \item What is the probability that one of the entries is exactly 16?
        Given A is the event that one entry is exactly 16. I'd break this question up into 3 parts.
        \begin{enumerate}
            \item The number of ways to choose 16 rolls out of 30: $\binom{30}{16}$
            \item The chance of one number 16 times: $ \frac{1}{6}^{16} $
            \item The change of getting every other number 14 times (the rest of the rolls): $ \frac{5}{6}^{14} $
            \item The final factor is the number of  
        \end{enumerate}
            \textbf{Solution:} 
        \[ P(A) = \binom{30}{16} \cdot \frac{1}{6}^{16} \cdot  \frac{5}{6}^{14} \]
        \item Explain why this is not a uniform random sampling. (Show that the probability of each outcome is not equally likely.)
        The following problem is not a uniform random sample because the outcomes of two elements have different probabilities so sampling can not be uniform.
        For example: 
        \begin{enumerate}
            \item Given event A = (30, 0, 0, 0, 0, 0)
            \item Given event B = ( 5, 5, 5, 5, 5, 5)
        \end{enumerate}
        The probability of A, $ P(A) = \frac{1}{6}^{30}$
        The probability of B, $ \frac{30!}{5!^{6}} \cdot \frac{1}{6}^{30}$
        \[ P(B) = \frac{30!}{5!^{6}} \cdot P(A)\]
        \[ P(B) > P(A) \]    

        The probability of event B happening is greater than event A thus the sampling is not uniform.

        \item Which elements are the least likely to occur and why?
        The elements least likely to occur are the ones where only one entry was rolled 30 times. There is only one possible way to roll a 1 or 2 30 times in a row. Other elements can be made with multiple sequences

    \end{enumerate}
    \item Suppose you are at a casino and there are three slot machines. You know that two slot
    machines each win with probability 0.5, and the other wins with a probability of 0.9.
    You pick one of the slot machines (equally at random) and you play 8 games and win them all.
    \begin{enumerate}
        \item What is the probability that you are at the slot machine that wins with probability of 0.9?
        To find the probability that you were using the machine with a winning percentage of 0.9 given you won 8 straight games we will define the following:
        \begin{enumerate}
            \item Event N is you picked 0.9 slot machine
            \item Event W is you won 8 straight games
        \end{enumerate}
        Now we need to find the probability $P(N|W)$. Using Bayes Theorem we can rewrite the equation into values that we know
        \[ P(N|W) = \frac{P(W|N) P(N)}{P(W)} \]    
        \[ P(N|W) = \frac{(0.9)^8 \cdot \frac{1}{3}}{\frac{1}{3} \cdot (0.5^8 + 0.5^8 + 0.9^8)} \]
        The $\frac{1}{3}$ cancels, and $(0.9)^8 = 0.43046721$ and $(0.5)^8 = 0.00390625$, so
        \[
        P(N|W) = \frac{0.43046721}{2(0.00390625) + 0.43046721} \approx 0.982
        \]

        \item  If you decide to stay at the same slot machine (after winning 8 games), what is the probability that your next game results in a win?
        The probability that you win your next game depends upon the machine you are at as well as the probability that you are at the machine after winning 8 straight games.
        Let event $A$ be that you win your next game. Then
        \[P(A) = P(N|W)(0.9) + (1-P(N|W))(0.5)\]
        Using $P(N|W) \approx 0.982$,
        \[
        P(A) \approx (0.982)(0.9) + (1-0.982)(0.5) \approx 0.893
        \]

    \end{enumerate}

    \item Partnership Dominos is a game that involves 4 players and a Double-6 Domino set.
    A Double-6 Domino set is a set of 28 different tiles such that each tile has two values, one on each side. The values are given by a number of pips or dots to signify the value. There is one domino for each combination of values from 0 through 6 (including doubles).
    At the start of the round, each of the 4 players (North, South, East, West) gets 7 tiles.
    \begin{enumerate}
        \item What is the probability that North has the double-six tile (the tile with 6 pips on either side)?

        There is 1 double-six tile out of 28, and North gets 7 tiles.
        \[
        \frac{1}{28} \cdot 7 = \frac{7}{28} = \frac{1}{4} = 0.25
        \]

        \item What is the probability that North has at least one double? (there are 7 doubles)

        There are 7 doubles and 21 non-doubles.
        \[
        P(\text{no doubles}) = \left(\frac{21}{28}\right)^{7} \approx 0.133
        \]
        \[
        P(\text{at least 1 double}) = 1 - \left(\frac{21}{28}\right)^{7} \approx 0.867
        \]

        \item What is the probability that North has at least one domino with a 6 and at least one domino with a blank?

        There are 7 dominos with a 6 and 7 dominos with a blank.
        \[
        \binom{21}{7} = 116280
        \]
        is the number of hands with no 6, and also the number of hands with no blank.
        Hands with neither a 6 nor a blank:
        \[
        \binom{15}{7} = 6435
        \]
        Total hands:
        \[
        \binom{28}{7} = 1184040
        \]
        Using inclusion-exclusion,
        \[
        1184040 - (116280 + 116280) + 6435 = 957915
        \]
        \[
        P = \frac{957915}{1184040} \approx 0.809
        \]

        \item What is the probability that North and South together hold all 7 tiles containing the value 6?

        North and South together have 14 tiles, so the total number of hands they can hold is
        \[
        \binom{28}{14} = 40116600
        \]
        After taking all 7 tiles with a 6, there are $28 - 7 = 21$ tiles left, and we still need 7 more tiles.
        \[
        \binom{21}{7} = 116280
        \]
        \[
        P = \frac{116280}{40116600} \approx 0.00290
        \]

        \item A ``Bad hand'' is a set of 7 dominos with at least four doubles. What is the probability that one of the players has a bad hand?

        There are 7 doubles. The ways for one player to have at least 4 doubles are
        \[
        \binom{7}{4}\binom{21}{3} + \binom{7}{5}\binom{21}{2} + \binom{7}{6}\binom{21}{1} + \binom{7}{7}\binom{21}{0}
        \]
        \[
        = 46550 + 4410 + 147 + 1 = 51108
        \]
        Total hands for one player:
        \[
        \binom{28}{7} = 1184040
        \]
        \[
        \frac{51108}{1184040} \approx 0.0432
        \]
        For any of the 4 players,
        \[
        0.0432 \cdot 4 \approx 0.173
        \]
    \end{enumerate}

    \item Fantasy Five is a lottery game you can play for \$1. The way it works is, from the set $\{1, 2, 3, \ldots, 39\}$, you select a subset of 5 different numbers.
    The lottery commissioner also selects 5 winning numbers uniformly at random.
    You win money based on how many of your numbers match with the winning numbers.
    Here are the payouts:
    \begin{itemize}
        \item 5 out of 5 pays \$67000
        \item 4 out of 5 pays \$401
        \item 3 out of 5 pays \$15
        \item 2 out of 5 pays \$1
        \item 1 out of 5 pays \$0
        \item 0 out of 5 pays \$0
    \end{itemize}
    \begin{enumerate}
        \item What is the expected number of matches you make with the winning numbers?

        You pick 5 numbers, there are 5 winning numbers, and 39 numbers total.
        \[
        P(\text{a number matches}) = \frac{5}{39} \approx 0.128
        \]
        Using a binomial distribution,
        \[
        E(\text{matches}) = 0.128 \cdot 5 \approx 0.641
        \]

        \item What is the expected profit you make from playing this game?

        The game costs \$1. The number of ways to choose 5 numbers is
        \[
        \binom{39}{5} = 575757
        \]
        \[
        P(\text{match 5}) = \frac{\binom{5}{5}\binom{34}{0}}{\binom{39}{5}} = \frac{1}{575757}
        \]
        \[
        P(\text{match 4}) = \frac{\binom{5}{4}\binom{34}{1}}{575757} = \frac{170}{575757}
        \]
        \[
        P(\text{match 3}) = \frac{\binom{5}{3}\binom{34}{2}}{575757} = \frac{5610}{575757}
        \]
        \[
        P(\text{match 2}) = \frac{\binom{5}{2}\binom{34}{3}}{575757} = \frac{59840}{575757}
        \]
        Expected profit is the payout from 5, 4, 3, and 2 matches, minus \$1.
        \[
        67000\left(\frac{1}{575757}\right) + 401\left(\frac{170}{575757}\right) + 15\left(\frac{5610}{575757}\right) + \frac{59840}{575757} - 1 \approx -\$0.52
        \]

        \item What is the expected number of times you must play this game before and including your first win of at least one dollar?

        Winning at least \$1 means matching 2, 3, 4, or 5 numbers.
        \[
        \frac{59840}{575757} + \frac{5610}{575757} + \frac{170}{575757} + \frac{1}{575757} = \frac{65621}{575757} \approx 0.114
        \]
        \[
        E(\text{games}) = \frac{1}{0.114} \approx 8.77
        \]
    \end{enumerate}

\end{enumerate}

\end{document}